% Part: normal-modal-logic
% Chapter: syntax-and-semantics
% Section: normal-modal-logics

\documentclass[../../../include/open-logic-section]{subfiles}

\begin{document}

\olfileid{nml}{syn}{nor}

\olsection{Logika Modal Normal}

Logika modal bisa dipandang minangka himpunan !!{formula}s
kang nyukupi aturan-aturan tartamtu. Ana sawetara
logika modal kang mligi narik kawigatèn amarga
interpretasi lan panganggoné, lan luwih-luwih amarga
landhesan téoriné wis dimangertèni kanthi apik.
Logika iki diarani logika modal normal.

\begin{defn}\ollabel{defn:modal-logic}
\emph{Logika modal}~$\Log L$ yaiku himpunan
!!{formula}s saka basa modal dhasar kang
\begin{enumerate}
\item ngemot kabèh tautologi proposisional,
\item katutup tumrap substitusi seragam, lan
\item katutup tumrap modus ponens, yaiku yèn $!A$
  lan $(!A \lif !B) \in \Log L$, mula
  $!B \in \Log L$.
\end{enumerate}
\end{defn}

\begin{ex}
Sawetara himpunan !!{formula}s kang nyukupi definisi iki,
sanajan ora pati narik kawigatèn, mula kalebu logika modal yaiku:
\begin{enumerate}
\item himpunan kabèh tautologi;
\item himpunan kabèh !!{formula}s (logika modal kang ora konsisten).
\end{enumerate}
\end{ex}

\begin{defn}\label{defn:normal-modal-logic}
Logika modal~$\Log L$ diarani \emph{normal} yèn lan mung yèn
\begin{enumerate}
\item ngemot kaloro formula iki:
\begin{align}
\tag{K} \Box(p \lif q) & \lif (\Box p \lif \Box q) \\
\tag{Dual} \Diamond p & \liff \lnot \Box \lnot p
\end{align}
\item lan katutup tumrap aturan peniscayaan, yaiku
  yèn $!A \in \Log L$, mula $\Box !A \in \Log L$.
\end{enumerate}
\end{defn}

Kita bakal nemoni rong cara utama netepaké logika modal.
Cara kapisan adhedhasar téori bukti: minangka himpunan
!!{formula}s kang bisa diturunké ing sistem
!!{derivation} tartamtu. Cara kapindho adhedhasar
téori modèl: minangka himpunan !!{formula}s kang valid
ing kelas modèl tartamtu. Iki padha karo kahanan
ing logika proposisional lumrah (utawa logika tataran
kapisan). Ing kono sawijining himpunan !!{formula}s
uga bisa ditemtokaké kanthi rong cara: minangka
himpunan tautologi kang bener ing saben pangaji nilai
bebener, utawa minangka himpunan kabèh !!{derivable}
!!{formula}s ing sistem !!{derivation} tartamtu.
Kanggo logika modal normal, rong cara iki bisa
ditrapaké kanthi modèl relasional lan sistem bukti
aksiomatik jinis Hilbert.
\end{document}
